%! TEX root = ./doc.tex

\documentclass[12pt, a4paper]{article}

\usepackage{amsmath}
\usepackage{amssymb}
%\usepackage{graphicx}
\usepackage[utf8]{inputenc}
\usepackage[T1,T2A]{fontenc}
\usepackage[bulgarian,english]{babel}
\usepackage{hyperref}
\usepackage{pgfplots}
\usepackage{csvsimple}
\pgfplotsset{compat=1.18}

\usepackage{tabularx}
\usepackage{algorithm}
\usepackage{algpseudocode}
\usepackage{esvect}

\usepackage{numprint}
\usepackage{microtype}
%\usepackage{biblatex}[backend=bibtex, style=authoryear]
%\selectlanguage{bulgarian}

\setcounter{secnumdepth}{1}

%% `Elsevier LaTeX' style
\bibliographystyle{elsarticle-num}
%%%%%%%%%%%%%%%%%%%%%%%
\title{Паралелизация на N-Body симулация върху графичен процесор}
\author{Борис Василев \\ \href{mailto:borisnv@uni-sofia.bg}{borisnv@uni-sofia.bg}}
\date{\small{Софийски Университет "Св. Климент Охридски" \\
	Факултет по Математика и Информатика
	} \\
	10.06.2025
}

\begin{document}

\maketitle

%% Group authors per affiliation:
\begin{sloppypar}

\noindent\makebox[\linewidth]{\rule{\textwidth}{0.4pt}}
\section*{Абстракт}
	Проблемът N-Body е класически проблем в компютърната физика, който изисква изчисляване на гравитационните взаимодействия между N тела. Този доклад разглежда паралелизацията на този проблем върху графичен процесор (GPU) с цел ускоряване на изчисленията. Използван OpenGL за имплементация на алгоритъма, като се сравняват резултатите с последователно изпълнение върху CPU. Резултатите показват значително ускорение при използване на GPU, особено при големи стойности на N.

\noindent\makebox[\linewidth]{\rule{\textwidth}{0.4pt}}

\section{Анализ на проблема}

Проблемът N-Body се състои в изчисляване на гравитационните сили между N тела, които могат да бъдат планети, звезди или други астрономически обекти. Всеки обект влияе на другите чрез гравитационна сила, която зависи от масата и разстоянието между тях. Целта е да се изчисли движението на телата във времето, което изисква многократни изчисления на взаимодействията.

Проблемът и методи за неговата паралелизация се разглеждат от самото създаване на машини, способни на паралелни изчисления \cite{aarseth2003gravitational}. Решения на проблема с N тела се използват във изчислителната астрофизика \cite{lacey1994merger} и в оценката на графичен хардуер \cite{playne2009benchmarking}. В литературата са се утвърдили няколко различни алгоритъма.

\subsection{Изчисление по двойки (All-Pairs Method)}

Най-простият метод за симулиране на гравитационните сили е пресмятането на ускоренията, породени от взаимодействията между всеки две тела. Алгоритъмът е разгледан в \cite{burtscher2011efficient,nylons2007fast,Grimm_2022}, а в \cite{spurzem1999direct} e разгледано негово разширение при разглеждане на негравитационни сили. Започвайки от закона за гравитацията на Нютон, можем да изразим ускорението на всяко тяло като сума от ускоренията, породени от всички останали тела. Нека като начало означим позициите на $N$-те тела с $\vv{r}_i$, масите им с $m_i$, разстоянието между телата $i$ и $j$ с $\vv{r}_{ij} = \vv{r}_i - \vv{r}_j$ и скоростите с $\vv{v}_i$. Тогава ускорението $\vv{a}_i$ може да се изрази като:

\begin{equation}
\begin{aligned}
\vv{a}_i &= \frac{\sum_{j \neq i} \vv{F}_{ij}}{m_i} \\
		  &= \frac{\sum_{j \neq i} \frac{G m_i m_j}{|\vv{r}_{ij}|^2} \frac{\vv{r}_{ij}}{|\vv{r}_{ij}|}}{m_i} \\
		  &= \sum_{j \neq i} \frac{G m_j}{|\vv{r}_{ij}|^3} \vv{r}_{ij}
\end{aligned}
\end{equation}

Така ускорението е приведено във лесна за пресмятане форма като не изисква нормализиране на вектора $\vv{r}_{ij}$, а само изчисляване на неговата дължина. Извършването на пълната симулация след това може да се извърши чрез Ойлерово интегриране на ускорението. Означаваме данните на $t$-та стъпка със горен индекс $(t)$. Освен това, нека $\Delta t$ е интеграционната стъпка. Забелязваме, че $\vv{F}_{ij} = \vv{F}_{ji}$ за всеки $i \neq j$. Така може да спестим пресмятане на ускоренията. Така псевдокода за цялостната симулация (за една стъпка) е описан в Алгоритъм \ref{alg:nbody}.

\begin{algorithm}
\caption{N-Body симулация по двойки обекти}
\label{alg:nbody}
\begin{algorithmic}[1]
\For{$i \gets 1$ \textbf{to} $N$}
	\State $\vv{v}_i^{(t+1)} \gets \vv{v}_i^{(t)}$
	\For{$j \gets 1$ \textbf{to} $i-1$}
		\State $\vv{diff} \gets \vv{r}_j^{(t)} - \vv{r}_i^{(t)}$
		\State $\vv{F} \gets G \cdot \frac{\vv{diff}}{|\vv{diff}|^3} \cdot \Delta t$
		\State $\vv{v}_i^{(t+1)} \gets \vv{v}_i^{(t+1)} + m_j \cdot \vv{F}$
		\State $\vv{v}_j^{(t+1)} \gets \vv{v}_j^{(t+1)} - m_i \cdot \vv{F}$
	\EndFor
\EndFor
\For{$i \gets 1$ \textbf{to} $N$}
	\State $\vv{r}_i^{(t+1)} \gets \vv{r}_i^{(t)} + \vv{v}_i^{(t+1)} \cdot \Delta t$
\EndFor
\end{algorithmic}
\end{algorithm}

Това е и алгоритъма, който се стремим да паралелизираме в този доклад. Той е лесен за разбиране и имплементиране, но има и недостатък: извършва $O(N^2)$ пресмятания, което го прави неефективен при голям брой тела.

\subsection{Йерархични алгоритми}

По-ефективни алгоритми за симулация на N-Body проблемите използват йерархични структури от данни, които позволяват да се избегне изчисляването на взаимодействия между всички двойки тела. Най-често се използва структура, разделяща пространството рекурсивно до изпълняване на определен критерий (например клетката съдържа $\leq k$ тела). Взаимодействията в клетка-листо се извършват по метода на всички двойки, а между клетки се използват приближения на 'далечното векторно поле' (far-field approximation). Популярни алгоритми са Barnes-Hut \cite{barnes1986hierarchical} и Fast Multipole Method \cite{greengard1987fast,dehnen2014fast}. Те са значително по-ефективни от метода на всички двойки, но също така са по-сложни за имплементиране и излизат от рамките на курсовата работа. В този доклад ще се фокусираме върху метода на всички двойки, тъй като той е най-прост и лесен за паралелизиране.

\section{Реализация}

\subsection{Проучване}
Реализация на N-Body симулация по двойки обекти върху графичен процесор (GPU) е разглеждана в дълбочина от различни автори. \cite{burtscher2011efficient,nylons2007fast,belleman2008high} разглеждат основно CUDA само върху хардуер на NVIDIA, докато \cite{Grimm_2022} поддържа и HIP върху AMD хардуер. Поради наличния хардуер, решихме да използваме OpenGL за имплементация на алгоритъма. OpenGL е широко използван стандарт за графични изчисления и е съвместим с повечето графични карти, независимо от производителя. Това позволява по-широка съвместимост и лесно разпространение на приложението. По този начин постигаме добри ускорения дори и върху вградени графични процесори на Intel. 

\subsection{Паралелизация}

Реализацията на алгоритъма до голяма степен се доближава до този в \cite{nylons2007fast}. Основната разлика е, че използваме OpenGl Compute с GLSL (OpenGL Shading Language) за изпълнение на изчисленията върху GPU, вместо CUDA. В \ref{alg:nbody} можем да видим, че for-циклите $@1...@9$ и $@11...@13$ трябва да се извършат последователно, а всеки от тях може да бъде паралелизиран. Иначе казано, съставяме два етапа на изчисление: първият етап изчислява ускоренията и скоростите на телата, а вторият етап обновява позициите им. Всеки от тези етапи може да бъде изпълнен паралелно за всяко тяло. Поради високия хардуерен паралелизъм на графичните процесори, може за $N$ тела да създадем $N$ паралелно работещи нишки. При този подход, обаче, на редове @6 и @7 получаваме конкурентен достъп до скоростите на телата, което изисква допълнителна синхронизация между нишките. Всякаква глобална (между произволни две нишки) синхронизация би довела до голямо забавяне върху графичен процесор поради самата му архитектура. За да избегнем това, може да 'платим' като не се възползваме от факта, че $\vv{F}_{ij} = \vv{F}_{ji}$ и да изчисляваме ускоренията не телата напълно отделно. Накрая, за да избегнем числени грешки, породени от твърде близки разстояния между телата, не позволяваме изчисленото разстояние да е по-малко от предварително избрана константа $\varepsilon$. Така получаваме алгоритъмът \ref{alg:nbody-parallel}. В \cite{nylons2007fast} се прави и опит за подобряване на производителността за малки стойности на $N$, но тук се стремим да постигнем максимално ускорение при големи стойности на $N$ и за това не разглеждаме тази оптимизация.

\begin{algorithm}
\caption{Паралелна N-Body симулация по двойки обекти}
\label{alg:nbody-parallel}
\begin{algorithmic}[1]

\For{$i \gets 1$ \textbf{to} $N$ in parallel}
	\State $\vv{v}_i^{(t+1)} \gets \vv{v}_i^{(t)}$
	\For{$j \gets 1$ \textbf{to} $N$}
		\If{$i = j$}
			\State \textbf{continue}
		\EndIf
		\State $\vv{diff} \gets \vv{r}_j^{(t)} - \vv{r}_i^{(t)}$
		\State $d \gets max(|\vv{diff}|, \varepsilon)$
		\State $\vv{F} \gets G \cdot \frac{\vv{diff}}{d^3} \cdot \Delta t$
		\State $\vv{v}_i^{(t+1)} \gets \vv{v}_i^{(t+1)} + m_j \cdot \vv{F}$
	\EndFor
\EndFor
\For{$i \gets 1$ \textbf{to} $N$ in parallel}
	\State $\vv{r}_i^{(t+1)} \gets \vv{r}_i^{(t)} + \vv{v}_i^{(t+1)} \cdot \Delta t$
\EndFor


\end{algorithmic}
\end{algorithm}

\subsection{Имплементация}

\subsubsection{GPU имплементация}
За реализиране на акгоритъма създаваме два OpenGL Compute Shader-а, които изпълняват двете стъпки в алгоритъм \ref{alg:nbody-parallel}. Телата се съхраняват в буфер във традиционна 'Array of Structs' (AoS) форма. Този буфер се закача във вид на шейдърна споделена памет (Shader Storage Buffer Object - SSBO) за да може елементите да се достъпват от всички нишки. В алгоритъма не се налага различни нишки да достъпват същия елемент от буфера, така че не се налага синхронизация между нишките. 

С оглед на архитектурата на графичния процесор, проблемът трябва да бъде разделен на блокове от нишки, представени в двумерна решетка. Поради споделения кеш между нишките в един Thread Block, е по-ефективно нишките в един блок да обработват тела, намиращи се в последователни клетки в паметта. Така при размер на блока (фиксиран за модел на графичния процесор) $p \times p$ и брой тела $N$, броят на блоковете е $\lceil N / p^2 \rceil$. При това разделение, всяка нишка изчислява ускорението на $p$ тела, като използва данните от всички останали тела. Така цялото изчисление е разпределено в решетка от нишки $\lceil N / p^2 \rceil p \times p$. Ясно е, че при закръглянето на броя блокове, някои нишки ще останат без работа, но това е неизбежно.

\subsubsection{CPU имплементация}

За сравнение, имплементацията на CPU е направена с помощта на C++ и реализира последователния алгоритъм \ref{alg:nbody} при същата структура на данните. Имплементацията използва и SIMD където е възможно чрез библиотеката за векторно смятане GLM \cite{riccio2010glm}.

\section{Тестовe}

\subsection{Хардуер}

Тестовете са проведени върху машина с процесор Intel Core i7-6700 и графичен процесор AMD Radeon RX 7600 с 2048 поточни процесора и операционна система EndeavourOS Linux 6.13.7. 

\subsection{Тестови данни}

Тестовите данни са генерирани случайно и са нормално разпределени с дисперсия $30$. За по-висока визуална атрактивност в точката $[0,0]$ e поставено тяло, чиято маса е сумата на останалите тела. Това позволява да се наблюдава как останалите тела се движат около него. За целите на тестовете са използвани стойности на $N$ от $10$ до $10^5$, като за всяка стойност са извършени поне 10 симулации с различни начални условия ($100$ за $N \leq 1000$).

\subsection{Резултати и анализ}

На фигура \ref{fig:time} е показано сравнение на времето за изпълнение на симулацията върху CPU и GPU. Времето е измерено в микросекунди ($\mu s$) и е логаритмично скалирано спрямо броя на телата $N$. От графиката се вижда, че времето за изпълнение на симулацията върху GPU е значително по-ниско от това върху CPU, особено при големи стойности на $N$. На фигура \ref{fig:speedup} е показано ускорението на GPU спрямо CPU, което достига до над $400$ пъти при $N = 10^5$. Това показва, че паралелизацията на N-Body симулацията върху GPU е изключително ефективна и може да се използва за значително ускоряване на изчисленията.

Остава да оценим оптималността на паралелизацията. Знаем, че имплементацията на последователния алгоритъм спестява около половината операции чрез съображението по-горе. Освен това, всяко ядро на процесора е способно да работи на честота до 4 GHz, докато използвания графичния процесор не може да поддържа честоти над 2.3 GHz за дълго време. Така знаейки и максималния хардуерен паралелизъм на графичния ускорител - $2048$ можем да изчислим максималното теоретично ускорение, което може да се постигне. То е:
\begin{equation}
	\begin{aligned}
		S_{max} &= 2048 \cdot \frac{2.3}{4} \cdot \frac{1}{2} \\
		&\approx 588.8
	\end{aligned}
\end{equation}

Това е теоретичното максимално ускорение, което може да се постигне при идеални условия. Виждаме, че реалното ускорение достига до $425$, и получаваме ефективност на паралелизацията $E = \frac{425}{588.8} \approx 0.722$, с което считаме, че паралелизацията е близкo до оптимална. Пълна таблица с резултати от тестовете е показана на фигура \ref{fig:table}.




\begin{figure}
\caption{Графика на времето за изпълнение на симулацията}
\label{fig:time}
\centering
\begin{tikzpicture}
\begin{axis}[
	%title={Сравнение на времето за изпълнение върху CPU и GPU},
	xlabel={log(N)},
	ylabel={log(Време) (us)},
	legend style={at={(0.02,0.75)}, anchor=south west},
	grid=major,
	width=0.8\textwidth,
	height=0.5\textwidth,
	xmode=log,
	ymode=log,
]
\addplot table [
	col sep=comma,
	x={N},
	y={CPU},
	] {./data.csv};
\addplot table [
	col sep=comma,
	x={N},
	y={GPU},
	] {./data.csv};
\legend{CPU, GPU};
\end{axis}
\end{tikzpicture}
\end{figure}

\begin{figure}
\caption{Графика ускорението на GPU спрямо CPU}
\label{fig:speedup}
\centering
\begin{tikzpicture}
\begin{axis}[
	title={},
	xlabel={log(N)},
	ylabel={Ускорение (CPU / GPU)},
	legend pos=outer north east,
	grid=major,
	width=0.8\textwidth,
	height=0.5\textwidth,
]
\addplot table [
	col sep=comma,
	x={N},
	y={Speedup},
	] {./data.csv};
\end{axis}
\end{tikzpicture}
\end{figure}

\npdecimalsign{.}
\nprounddigits{2}

\begin{figure}
\caption{Таблица с резултати от тестовете}
\label{fig:table}
\centering
\begin{tabularx}{.95\textwidth}{|X|X|X|n{4}{2}|}%
	\hline
	\bfseries N & \bfseries {CPU time ($\mu s$)} & \bfseries {GPU time ($\mu s$)} & \bfseries Speedup
\csvreader[head to column names]{data.csv}{}
	{\\\hline \N & \CPU & \GPU & \Speedup} 
	\\
	\hline
\end{tabularx}
\end{figure}

\end{sloppypar}
\clearpage

\bibliography{refs}

\end{document}
